\chapter*{前言与阅读约定}
\addcontentsline{toc}{chapter}{前言与阅读约定}
本书以吉田耕作《泛函分析》第六版中译本为主线，以英文第六版辅助核对术语、目录和页码。第一部分保留原书第零章至第十四章的章次与各节主题，包括第五章和第十章的原书附录。目前第零至第十一章共 110/132 个编号节已有首轮正文，附 186 道解答习题；第五章与第十章附录已写，第十二至十四章仍保留占位。已有正文的比例不代表阅读进度、工作量比例或最终审校完成率。

主线从预备知识、半范数与 Banach 空间基本定理进入 Hilbert 空间、Hahn--Banach、强弱收敛及对偶算子，再连接 Fourier 与微分方程、预解式和谱、半群及发展方程。第八章加入两类平均过程、殆周期均值与全纯演算；第九章展开生成元、Hille--Yosida、Lumer--Phillips、Stone、全纯半群、分数幂、Trotter--Kato 及对偶连续部分。第十章连接紧性判别、核算子、Sobolev 紧嵌入、Fredholm 谱论与 Dirichlet 弱解，并补写核空间附录。第十一章从赋范代数进入谱测度、无界自伴演算、紧群表示和 Wiener 逼近。每章核查记录区分正文覆盖、作者自查、排版验收与独立数学审校。

部分一般结论仍作为外部前提或待展开内容，例如 Riesz--Markov、一般 Eberlein--\v Smulian、Titchmarsh 与变系数正则性，以及部分局部凸推广、扇形预解式的反向构造和一般从属核。具体假设、直接查阅范围与证明边界在各节及 COVERAGE 中注明；不把首轮写成或编译通过等同于独立数学审校完成。

第二部分“广义函数专题扩展”接入此前已经写成的八章内容，章号顺接为第十五章至第二十二章。它以施瓦尔茨《广义函数论》和盖尔芳特、希烙夫《广义函数 III：微分方程理论的若干问题》为专题参考。吉田原书本身涉及广义函数的各节仍保留在第一部分，后续写作可引用专题中的详细证明，也可补充原书特有的论述。本次调整保持现有专题正文与习题，采用标签引用适应新的章号。

目录框架中的章节名与节名以中译本为准，英文版辅助核对人名与对应主题。每个主线占位节记录两版印刷起始页；附录另列 PDF 页序的换算。后续应逐节补充定义、定理、证明、例子和习题，并统一记号。下面的配对与 Fourier 约定目前适用于已写成的广义函数分部，主线正文写作时也应明确它与原书约定的差别。

\Needspace{10\baselineskip}
\section*{符号与配对}
\phantomsection
\addcontentsline{toc}{section}{符号与配对}
始终令 $\Omega\subset\R^n$ 为开集，$K\Subset\Omega$ 表示紧集 $K$ 包含于 $\Omega$。多重指标 $\alpha=(\alpha_1,\ldots,\alpha_n)\in\N_0^n$，并约定
\[
 |\alpha|=\sum_j\alpha_j,\quad \alpha!=\prod_j\alpha_j!,\quad
 x^\alpha=\prod_jx_j^{\alpha_j},\quad
 \partial^\alpha=\partial_{x_1}^{\alpha_1}\cdots\partial_{x_n}^{\alpha_n}.
\]
分布是复线性泛函，$\Pair{T}{\varphi}$ 对 $\varphi$ 线性，正则分布采用
\[
 \Pair{T_f}{\varphi}=\int_\Omega f(x)\varphi(x)\,\dd x,
\]
不在测试函数上加复共轭。Hilbert 空间的内积则另行采用 $ (u,v)=\int u\overline v$，对第一变量线性。这两种约定不可混同。

\begin{center}
\begin{tabular}{ll}
\toprule
符号 & 含义\\\midrule
$\Test(\Omega)$ & $C_c^\infty(\Omega)$，带标准 LF 拓扑\\
$\Dist(\Omega)$ & $\Test(\Omega)$ 的连续线性对偶\\
$\Smooth(\Omega)$ & $C^\infty(\Omega)$，带紧集上各阶一致收敛拓扑\\
$\CompactDist(\Omega)$ & 紧支集分布空间\\
$\Sch(\R^n),\Temp(\R^n)$ & Schwartz 空间与温和分布空间\\
$\supp T$ & 分布 $T$ 的支集\\
$T_j\to T$ & 默认指逐测试函数收敛，即弱对偶收敛\\
\bottomrule
\end{tabular}
\end{center}

傅里叶变换统一使用角频率约定：
\begin{equation}\label{eq:fourier-convention}
 \widehat f(\xi)=\int_{\R^n}e^{-ix\cdot\xi}f(x)\,\dd x,
 \qquad f(x)=(2\pi)^{-n}\int_{\R^n}e^{ix\cdot\xi}\widehat f(\xi)\,\dd\xi.
\end{equation}
施瓦尔茨原书部分公式使用 $2\pi$ 频率约定，盖尔芳特、希烙夫原书相关章节采用正号指数；查阅时须作符号与常数换算。

\section*{建议的阅读路径}
主线按吉田原书从预备知识读至发展方程。当前可直接阅读的专题正文从第 \ref{ch:motivation} 章起：第 \ref{ch:test}--\ref{ch:operations} 章建立测试空间与分布运算，第 \ref{ch:convolution}--\ref{ch:fourier} 章讨论卷积与 Fourier 变换，第 \ref{ch:pde} 章讨论微分方程模型，第 \ref{ch:weighted} 章介绍加权空间与广义固有函数。已写成的习题解答集中置于书后；新增占位节目前没有虚构的习题或证明。
